\documentclass[10pt,a4paper]{amsart}
\usepackage[margin=19mm]{geometry}
\usepackage{amsmath,amssymb,mathtools}
\usepackage[scr=rsfs,cal=euler]{mathalfa}
\usepackage{xfrac}
\usepackage[unicode,hidelinks,bookmarks=false]{hyperref}
\newcommand{\ie}{i.e.\ }
\newcommand{\cf}{cf.\ }
\newcommand{\resp}{respectively\ }
\newcommand{\Subsection}{\subsection}
\providecommand{\nobreakdash}{\nobreak\hbox{-}}
\setlength{\emergencystretch}{2em}
\multlinegap0pt
\usepackage{bm}
\usepackage{bbm}
\usepackage{dsfont}
\usepackage{xspace}
\usepackage{cancel}
\usepackage{mathtools}
\usepackage[shortlabels]{enumitem}
\usepackage{amsthm}
\makeatletter
\@ifundefined{theorem}{\newtheorem{theorem}{Theorem}}{}
\makeatother
\title[Injective (edge) colorings of generalized Sierpi\'{n}ski gra]{Injective (edge) colorings of generalized Sierpi\'{n}ski graphs}
\author{C. K. Bhanupriya, Bo\v{s}tjan Bre\v{s}ar}
\date{Source: arXiv:2508.14479v2 (CC BY 4.0)}
\begin{document}
\maketitle

\noindent\emph{Source and licence.} Injective (edge) colorings of generalized Sierpi\'{n}ski graphs. Version arXiv:2508.14479v2 (CC BY 4.0), distributed under the Creative Commons Attribution 4.0 International licence, as recorded in the arXiv licence metadata for this exact version and shown on the abstract page https://arxiv.org/abs/2508.14479v2. Full rights notice: licenses/arxiv-2508.14479-CC-BY-4.0.txt.\par\smallskip

\noindent\textit{Editorial scope.} Counted unit: the Theorem whose source label is \texttt{thm:triangle-free} in the article (source0.tex lines 1039--1049, source label \texttt{thm:triangle-free}), delivered together with the article's own complete original proof of it (source0.tex lines 1050--1068, one \texttt{proof} environment of 19 lines). No mathematics is written, repaired or re-derived: statement and proof are the article's own text line for line. Required context retained uncounted: none. Every definition, display and label of those lines is retained unchanged. Omitted from this package and not redistributed: 1--1038, 1069--1717. Nothing inside a retained excerpt is omitted or altered: every retained body is delivered exactly as the article's lines are written, so no omission receipt of any kind accompanies this package. Citations: the retained excerpts carry no citation command at all, so no bibliography entry of the article is transcribed and no reference is redistributed. Reference adaptation: none was needed and none was made; every reference inside the delivered unit resolves inside it, so no unresolved reference or citation is produced. Printed slips kept literal: no printed word, symbol, hypothesis or number of the counted statement or of its proof was altered. Layout and class: the article's own class is \texttt{article} with packages including graphicx, amsthm, amsfonts, amsbsy, amssymb, amsmath, graphics, babel, lineno, inputenc, enumerate, enumitem, tkz-berge, hyperref; the wrapper is the lane's pinned \texttt{amsart} editorial wrapper at 10pt on A4, which restates the theorem-like environments the retained text needs and adds the article's own macro definitions that the retained text needs (none), with extra wrapper packages (bm, bbm, dsfont, xspace, cancel, mathtools, enumitem). The delivered numbering is the unit's own. Assets: no retained span contains a figure environment, a graphics-inclusion command, a PDF-page inclusion, a table or a TikZ picture; no image file is copied into this package, the package declares no asset dependency and no image file is redistributed with it. Licence: version arXiv:2508.14479v2 of this article is distributed under the Creative Commons Attribution 4.0 International licence, as recorded in the arXiv licence metadata for this exact version and shown on the abstract page https://arxiv.org/abs/2508.14479v2. A full rights notice is delivered at licenses/arxiv-2508.14479-CC-BY-4.0.txt. Characters: every retained excerpt is pure ASCII, so the delivered engine, XeLaTeX with its default Latin Modern text font, reports no missing character.
\begin{theorem}
\label{thm:triangle-free}
For any triangle-free graph $G$ and any positive integer $n\ge 3$,
	$$\chi_i'(S_G^3)\le \chi_i'(S_G^n)\le \chi_i'(S_G^3)+1.$$
In addition, $\chi_i'(S_G^n)=\chi_i'(S_G^3)$ for all $n\ge 3$ if
\begin{enumerate}[label=(\roman*),noitemsep]
    \item there exists a $\chi_i'$-coloring of $S_G^3$, where $\chi_i'(S_G^3)=k$, such that every vertex in $S_G^3$ incident to an edge with color $k$ is at distance at least $3$ from every extreme vertex $(u,u,u)$, $u\in V(G)$, of $S_G^3$; or
    \item $G$ is bipartite with the bipartition $V(G)=A\cup B$ and there exists a $\chi_i'$-coloring of $S_G^3$, where $\chi_i'(S_G^3)=k$, such that for every $u\in A$ the edges incident with the extreme vertex $(u,u,u)$ receive color $k$, and every vertex in $S_G^3$ incident to an edge with color $k$ is at distance at least $2$ from every extreme vertex $(v,v,v)$, where $v\in B$.
\end{enumerate}
\label{thm:SG3}
\end{theorem}

\begin{proof}
Clearly, $\chi_i(S_G^n)\geq \chi_i(S_G^3)$, since $S_G^3$ is a subgraph of $S_G^n$.

Let $f:E(S_G^3)\rightarrow [k]$ be a $\chi_i'$-coloring of $S_G^3$. For any $n\ge 4$ we will introduce a coloring $f_1:E(S_G^n)\rightarrow [k+1]$ for which we will prove it is an injective edge coloring of $S_G^n$. Note that $V(S_G^n)$ can be partitioned into $|V(G)|^{n-3}$ subsets each of which induces the graph $S_G^3$.  First, we assign the color $k+1$ to the edges between different copies of $S_G^3$ in $S_G^n$. That is, for any edge $(u_1,\ldots,u_{n-3},v,v,v)(v_1,\ldots,v_{n-3},u,u,u)$, where $u_i,v_i\in V(G)$ for all $i\in [n-3]$ and $uv\in E(G)$, we assign $$f_1((u_1,\ldots,u_{n-3},v,v,v)(v_1,\ldots,v_{n-3},u,u,u))=k+1$$ 
(note that there exists $d\in [n-4]$ such that $u_i=v_i$ for all $i<d$, $u_d=u$, $v_d=v$, and $u_i=v$, $v_i=u$ for all $i\in\{d+1,\ldots,n-4\}$).
Concerning the edges whose endvertices belong to one and the same copy of $S_G^3$, we color them in the same way as they are colored by $f$ with some exceptions that we specify next. For this purpose, for each edge $uv\in E(G)$ we arbitrarily choose one endvertex, either $u$ or $v$, and we call it the chosen vertex of the edge $uv$. Now, consider a copy of $S_G^3$ in $S_G^n$ and let $(w_1,\ldots,w_{n-3})\in V(G)^{n-3}$ be the vertex to which it corresponds; that is, consider the subgraph $S_{(w_1,\ldots,w_{n-3})}$ of $S_G^n$ induced by $$\{(w_1,\ldots,w_{n-3},x,y,z):\, \textrm{ where }(x,y,z)\in V(G)^3\}.$$
If a vertex in $S_{(w_1,\ldots,w_{n-3})}$ is an extreme vertex of that copy of $S_G^3$, say a vertex $(w_1,\ldots,w_{n-3},x,x,x)$, which has a neighbor $(u_1,\ldots,u_{n-3},y,y,y)$ in another copy of $S_G^3$, namely $S_{(u_1,\ldots,u_{n-3})}$, and $x$ is the chosen endvertex of the edge $xy\in E(G)$, then $f_1$ colors all edges of $S_{(w_1,\ldots,w_{n-3})}$ incident with $(w_1,\ldots,w_{n-3},x,x,x)$ by color $k+1$. On the other hand, $f_1$ colors all other edges between vertices in $S_{(w_1,\ldots,w_{n-3})}$ with the same color as $f$ colors the corresponding edges in $S_G^3$. That is, $f_1(w_1,\ldots,w_{n-3},u,v,w)(w_1,\ldots,w_{n-3},x,y,z))=f((u,v,w)(x,y,z))$ for all such edges. 

To prove that $f_1$ is indeed an injective edge coloring of $S_G^n$, first consider the edges colored by color $k+1$. Note that the distance between any two extreme vertices in $S_G^3$ is at least $7$. Hence, if two edges that are in the same copy of $S_G^3$ receive the color $k+1$ by $f_1$, then either they have a common endvertex, which is an extreme vertex of that copy, or their distance is at least $5$ (since they are incident to distinct extreme vertices of that copy of $S_G^3$ in $S_G^n$). In the former case, they do not have a common edge, since $G$ is triangle-free. In the latter case, since their distance is at least $5$, they also do not have a common edge.  
Next, the distance between two edges colored with color $k+1$ that are from distinct copies of $S_G^3$ is at least $7$. Finally, the distance from an edge $e$ with color $k+1$ that lies between two distinct copies of $S_G^3$ in $S_G^n$ to any other edge with color $k+1$ is at least $6$, unless an edge with color $k+1$ is adjacent to $e$ (in which case they clearly do not have any common edges). 

Second, concerning the edges $e$ with $f_1(e)\le k$, we distinguish two cases. If two edges $e$ and $e'$ with $f_1(e)=f_1(e')$ belong to the same copy of $S_G^3$ in $S_G^n$, then since $f_1(e)=f(e)=f(e')=f_1(e')$, they do not have a common edge, since $f$ is an injective edge coloring of $S_G^3$. On the other hand, if $e$ and $e'$ are edges from distinct copies of $S_G^3$, and $f_1(e)=f_1(e')\le k$, then by the construction of $f_1$, they also cannot have a common edge. Indeed, on any path between $e$ and $e'$ there lie two edges colored with color $k+1$, hence $e$ and $e'$ do not have a common edge. Thus, $f_1$ is an injective edge coloring of $S_G^n$, which uses $k+1$ colors, so the upper bound $\chi_i'(S_G^n)\le \chi_i'(S_G^3)+1$ is proved. 

For the proof of {\it (i)}, assume that there exists a $\chi_i'$-coloring of $S_G^3$, where $\chi_i'(S_G^3)=k$, such that every vertex in $S_G^3$ incident to an edge with color $k$ is at distance at least $3$ from every extreme vertex $(u,u,u)$, $u\in V(G)$, of $S_G^3$. Then $f_1$ can be defined in almost the same way as in the previous (general) case, except that all edges that $f_1$ colored to $k+1$ in the previous case are recolored to $k$. Indeed, the condition that the distance from every extreme vertex $(u,u,u)$, $u\in V(G)$, of $S_G^3$, to a vertex incident with an edge with color $k$ is at least $3$ ensures that no two edges with color $k$ will have a common edge.

The proof of {\it (ii)} is straightforward. We define the coloring $f_1$ by 
$$f_1((w_1,\ldots,w_{n-3},u,v,w)(w_1,\ldots,w_{n-3},x,y,z))=f((u,v,w)(x,y,z))$$ for every edge in $S_G^n$ whose endvertices lie in the same copy of $S_G^3$ of $S_G^n$, where $f$ is a coloring of $S_G^3$ that satisfies the properties in the assumption of the statement. That is, in each copy of $S_G^3$ we use the coloring $f$. Now, to each edge between distinct copies of $S_G^3$ in $S_G^n$ let $f_1$ assign color $k$. Note that an edge between two vertices $(u_1,\ldots, u_n),(v_1,\ldots,v_n)\in V(S_G^n)$ that belong to distinct copies of $S_G^3$ exists only if $u_{n-2}=u_{n-1}=u_{n}=u$ and $v_{n-2}=v_{n-1}=v_{n}=v$ and $uv\in E(G)$. Since $G$ is bipartite, we may assume that $u\in A$ and $v\in B$. Note that all edges incident with $(u_1,\ldots, u_n)$ are colored with color $k$ by $f_1$, and since $G$ is triangle-free they do share a common edge. On the other hand, the distance from $(v_1,\ldots,v_n)$ to any vertex in the same copy of $S_G^3$, which is incident to an edge with color $k$, is at least $2$. Therefore, the edge $(u_1,\ldots, u_n)(v_1,\ldots,v_n)$ has no common edge to any edge with color $k$. 
The proof is complete.  
\end{proof}	

\end{document}
